% Options for packages loaded elsewhere
\PassOptionsToPackage{unicode}{hyperref}
\PassOptionsToPackage{hyphens}{url}
\PassOptionsToPackage{dvipsnames,svgnames,x11names}{xcolor}
\documentclass[
  english,
  11pt,
  a4paper,
]{article}
\usepackage{xcolor}
\usepackage[margin=2.4cm]{geometry}
\usepackage{amsmath,amssymb}
\setcounter{secnumdepth}{-\maxdimen} % remove section numbering
\usepackage{iftex}
\ifPDFTeX
  \usepackage[T1]{fontenc}
  \usepackage[utf8]{inputenc}
  \usepackage{textcomp} % provide euro and other symbols
\else % if luatex or xetex
  \usepackage{unicode-math} % this also loads fontspec
  \defaultfontfeatures{Scale=MatchLowercase}
  \defaultfontfeatures[\rmfamily]{Ligatures=TeX,Scale=1}
\fi
\usepackage{lmodern}
\ifPDFTeX\else
  % xetex/luatex font selection
  \setmainfont[]{Liberation Serif}
  \setmonofont[Scale=0.82]{DejaVu Sans Mono}
\fi
% Use upquote if available, for straight quotes in verbatim environments
\IfFileExists{upquote.sty}{\usepackage{upquote}}{}
\IfFileExists{microtype.sty}{% use microtype if available
  \usepackage[]{microtype}
  \UseMicrotypeSet[protrusion]{basicmath} % disable protrusion for tt fonts
}{}
\makeatletter
\@ifundefined{KOMAClassName}{% if non-KOMA class
  \IfFileExists{parskip.sty}{%
    \usepackage{parskip}
  }{% else
    \setlength{\parindent}{0pt}
    \setlength{\parskip}{6pt plus 2pt minus 1pt}}
}{% if KOMA class
  \KOMAoptions{parskip=half}}
\makeatother
\usepackage{color}
\usepackage{fancyvrb}
\newcommand{\VerbBar}{|}
\newcommand{\VERB}{\Verb[commandchars=\\\{\}]}
\DefineVerbatimEnvironment{Highlighting}{Verbatim}{commandchars=\\\{\}}
% Add ',fontsize=\small' for more characters per line
\newenvironment{Shaded}{}{}
\newcommand{\AlertTok}[1]{\textcolor[rgb]{1.00,0.00,0.00}{\textbf{#1}}}
\newcommand{\AnnotationTok}[1]{\textcolor[rgb]{0.38,0.63,0.69}{\textbf{\textit{#1}}}}
\newcommand{\AttributeTok}[1]{\textcolor[rgb]{0.49,0.56,0.16}{#1}}
\newcommand{\BaseNTok}[1]{\textcolor[rgb]{0.25,0.63,0.44}{#1}}
\newcommand{\BuiltInTok}[1]{\textcolor[rgb]{0.00,0.50,0.00}{#1}}
\newcommand{\CharTok}[1]{\textcolor[rgb]{0.25,0.44,0.63}{#1}}
\newcommand{\CommentTok}[1]{\textcolor[rgb]{0.38,0.63,0.69}{\textit{#1}}}
\newcommand{\CommentVarTok}[1]{\textcolor[rgb]{0.38,0.63,0.69}{\textbf{\textit{#1}}}}
\newcommand{\ConstantTok}[1]{\textcolor[rgb]{0.53,0.00,0.00}{#1}}
\newcommand{\ControlFlowTok}[1]{\textcolor[rgb]{0.00,0.44,0.13}{\textbf{#1}}}
\newcommand{\DataTypeTok}[1]{\textcolor[rgb]{0.56,0.13,0.00}{#1}}
\newcommand{\DecValTok}[1]{\textcolor[rgb]{0.25,0.63,0.44}{#1}}
\newcommand{\DocumentationTok}[1]{\textcolor[rgb]{0.73,0.13,0.13}{\textit{#1}}}
\newcommand{\ErrorTok}[1]{\textcolor[rgb]{1.00,0.00,0.00}{\textbf{#1}}}
\newcommand{\ExtensionTok}[1]{#1}
\newcommand{\FloatTok}[1]{\textcolor[rgb]{0.25,0.63,0.44}{#1}}
\newcommand{\FunctionTok}[1]{\textcolor[rgb]{0.02,0.16,0.49}{#1}}
\newcommand{\ImportTok}[1]{\textcolor[rgb]{0.00,0.50,0.00}{\textbf{#1}}}
\newcommand{\InformationTok}[1]{\textcolor[rgb]{0.38,0.63,0.69}{\textbf{\textit{#1}}}}
\newcommand{\KeywordTok}[1]{\textcolor[rgb]{0.00,0.44,0.13}{\textbf{#1}}}
\newcommand{\NormalTok}[1]{#1}
\newcommand{\OperatorTok}[1]{\textcolor[rgb]{0.40,0.40,0.40}{#1}}
\newcommand{\OtherTok}[1]{\textcolor[rgb]{0.00,0.44,0.13}{#1}}
\newcommand{\PreprocessorTok}[1]{\textcolor[rgb]{0.74,0.48,0.00}{#1}}
\newcommand{\RegionMarkerTok}[1]{#1}
\newcommand{\SpecialCharTok}[1]{\textcolor[rgb]{0.25,0.44,0.63}{#1}}
\newcommand{\SpecialStringTok}[1]{\textcolor[rgb]{0.73,0.40,0.53}{#1}}
\newcommand{\StringTok}[1]{\textcolor[rgb]{0.25,0.44,0.63}{#1}}
\newcommand{\VariableTok}[1]{\textcolor[rgb]{0.10,0.09,0.49}{#1}}
\newcommand{\VerbatimStringTok}[1]{\textcolor[rgb]{0.25,0.44,0.63}{#1}}
\newcommand{\WarningTok}[1]{\textcolor[rgb]{0.38,0.63,0.69}{\textbf{\textit{#1}}}}
\usepackage{longtable,booktabs,array}
\usepackage{calc} % for calculating minipage widths
% Correct order of tables after \paragraph or \subparagraph
\usepackage{etoolbox}
\makeatletter
\patchcmd\longtable{\par}{\if@noskipsec\mbox{}\fi\par}{}{}
\makeatother
% Allow footnotes in longtable head/foot
\IfFileExists{footnotehyper.sty}{\usepackage{footnotehyper}}{\usepackage{footnote}}
\makesavenoteenv{longtable}
\ifLuaTeX
\usepackage[bidi=basic,shorthands=off,]{babel}
\else
\usepackage[bidi=default,shorthands=off,]{babel}
\fi
\ifPDFTeX
\else
\babelfont{rm}[]{Liberation Serif}
\fi
\ifLuaTeX
  \usepackage{selnolig} % disable illegal ligatures
\fi
\setlength{\emergencystretch}{3em} % prevent overfull lines
\providecommand{\tightlist}{%
  \setlength{\itemsep}{0pt}\setlength{\parskip}{0pt}}
% Dryad Language Specification — LaTeX preamble
% Loaded with --include-in-header. Compatible with XeLaTeX/Tectonic.

\usepackage{graphicx}
\usepackage{fancyhdr}
\usepackage{titlesec}
\usepackage{enumitem}
\usepackage{microtype}
\usepackage{booktabs}
\usepackage{array}
\usepackage{xcolor}
\usepackage{needspace}

% ---------------------------------------------------------------- colours
\definecolor{dryadgreen}{RGB}{47,93,58}
\definecolor{dryadlight}{RGB}{110,140,115}
\definecolor{rulegray}{RGB}{213,219,212}

% ---------------------------------------------------------------- headings
\titleformat{\section}
  {\Large\bfseries\color{dryadgreen}}{}{0pt}{}
  [\vspace{-6pt}{\color{dryadgreen}\rule{\linewidth}{0.9pt}}]
\titleformat{\subsection}
  {\large\bfseries\color{dryadgreen!85!black}}{}{0pt}{}
\titleformat{\subsubsection}
  {\normalsize\bfseries\color{black}}{}{0pt}{}
\titleformat{\paragraph}[runin]
  {\normalsize\bfseries}{}{0pt}{}[.]
\titlespacing*{\section}{0pt}{20pt plus 4pt minus 2pt}{9pt}
\titlespacing*{\subsection}{0pt}{15pt plus 3pt minus 2pt}{6pt}
\titlespacing*{\subsubsection}{0pt}{12pt plus 2pt minus 2pt}{4pt}

% ---------------------------------------------------------------- layout
\setlength{\parindent}{0pt}
\setlength{\parskip}{0.55em}
\renewcommand{\arraystretch}{1.18}
\setlist{itemsep=0.15em, topsep=0.35em, parsep=0.2em}

% ---------------------------------------------------------------- headers
\pagestyle{fancy}
\fancyhf{}
\fancyhead[L]{\small\color{dryadgreen}\textsc{Dryad Language Specification}}
\fancyhead[R]{\small\color{dryadgreen}\textsc{Edition 1.0 --- Draft}}
\fancyfoot[C]{\small\thepage}
\renewcommand{\headrulewidth}{0.35pt}
\renewcommand{\footrulewidth}{0pt}
\fancypagestyle{plain}{%
  \fancyhf{}%
  \fancyfoot[C]{\small\thepage}%
  \renewcommand{\headrulewidth}{0pt}%
}

% ---------------------------------------------------------------- title page
\makeatletter
\renewcommand{\maketitle}{%
  \begin{titlepage}
    \thispagestyle{empty}
    \centering
    \vspace*{1.6cm}
    \includegraphics[width=5.6cm]{assets/dryad-logo.png}\par
    \vspace{1.3cm}
    {\color{dryadgreen}\rule{\linewidth}{1.1pt}}\par
    \vspace{0.7cm}
    {\fontsize{30}{36}\selectfont\bfseries\color{dryadgreen}\@title\par}
    \vspace{0.7cm}
    {\Large\color{dryadlight}\@author\par}
    \vspace{1.1cm}
    {\color{dryadgreen}\rule{\linewidth}{1.1pt}}\par
    \vfill
    {\large\bfseries Edition 1.0 --- Draft}\par
    \vspace{0.25cm}
    {\large \@date}\par
    \vspace{0.25cm}
    {\normalsize Dryad Language Project}\par
    \vspace{1.2cm}
    \begin{minipage}{0.86\linewidth}
      \small\itshape
      This document specifies the Dryad programming language. It is the
      normative reference for language implementors. Clauses are normative
      unless explicitly marked informative; proposals and implementation
      guidance are confined to the annexes.\par
      \vspace{0.5em}
      \upshape\small Licensed under the MIT License.
    \end{minipage}
  \end{titlepage}
  \setcounter{page}{2}%
}
\makeatother
\usepackage{bookmark}
\IfFileExists{xurl.sty}{\usepackage{xurl}}{} % add URL line breaks if available
\urlstyle{same}
\hypersetup{
  pdftitle={Dryad Language Specification},
  pdfauthor={Pedro Kaleb},
  pdflang={en},
  colorlinks=true,
  linkcolor={dryadgreen},
  filecolor={Maroon},
  citecolor={Blue},
  urlcolor={dryadgreen},
  pdfcreator={LaTeX via pandoc}}

\title{Dryad Language Specification}
\author{Pedro Kaleb}
\date{October 2026}

\begin{document}
\maketitle
\begin{abstract}
This document specifies the Dryad programming language: its lexical
grammar, type system, expressions, statements, declarations, module
system, execution model, error model, concurrency primitives, foreign
function interface, runtime architecture, standard library, bytecode
instruction set, and ahead-of-time compilation pipeline. It is the
normative reference for implementors of Dryad and supersedes the Dryad
Conventions 1.0 draft as the authoritative specification.
\end{abstract}

{
\hypersetup{linkcolor=black}
\setcounter{tocdepth}{3}
\tableofcontents
}
\section{Status of This Document}\label{status-of-this-document}

\textbf{Document identifier.} DRYAD-SPEC-1.0-DRAFT.

\textbf{Status.} Draft Standard. This edition is complete for review but
is not yet a ratified standard. Implementations claiming conformance to
this edition MUST identify the edition and the conformance class they
implement (Clause 2).

\textbf{Normative content.} Clauses 1 through 20 and Annexes A through D
are normative. The Introduction, Clause 21, Clause 24, and Annexes E
through G are informative. A clause or annex is normative unless it is
explicitly marked \emph{informative}. Within normative clauses, text
marked \textbf{NOTE} is informative.

\textbf{Relationship to prior documents.} This document consolidates and
replaces the content of \emph{Dryad Conventions 1.0} (the
\texttt{conventions} publication). It also supersedes \emph{Dryad
Theoretical Foundation v2.0} as the authoritative language
specification; the v2 document is retained as an informative
architecture paper. Where this document and any prior document disagree,
this document prevails.

\textbf{Editorial decisions.} The conventions draft contains internal
inconsistencies. Every deliberate resolution is recorded in Annex E. An
implementation that follows Annex E is following this specification, not
deviating from it.

\textbf{Change control.} Changes to normative content require a new
edition or a published erratum. Editorial changes that do not alter
normative behaviour may be published as a revision of the same edition.

\textbf{Licence.} This specification is published under the MIT Licence.

\section{Introduction}\label{introduction}

\emph{This Introduction is informative.}

Dryad is a multi-paradigm programming language designed for clarity,
portability, and deterministic execution. Its syntax is deliberately
close to JavaScript and TypeScript, while its runtime architecture
borrows the micro-kernel discipline of systems languages: a small native
core exposes a fixed set of intrinsics, and the entire standard library
is written in Dryad itself.

The language is defined by the interaction of four layers:

\begin{enumerate}
\def\labelenumi{\arabic{enumi}.}
\tightlist
\item
  \textbf{Source language.} A statically scoped, dynamically typed
  language with optional strict typing.
\item
  \textbf{Execution tiers.} A tree-walking interpreter, a stack-based
  bytecode virtual machine, and an ahead-of-time compiler. The tiers are
  required to be observationally equivalent.
\item
  \textbf{Native boundary.} A minimal, explicitly specified intrinsic
  interface through which all operating-system interaction occurs.
\item
  \textbf{Standard library.} Pure Dryad modules that define the
  user-visible platform, including I/O, networking, collections,
  concurrency, and diagnostics.
\end{enumerate}

This specification is written for implementors. It does not teach the
language, and it does not describe the current implementation. It
defines what a conforming implementation MUST do. Where the current
implementation differs, the implementation is non-conforming; Clause 24
and Annex E identify known gaps for planning purposes.

\subsection{Design goals}\label{design-goals}

\begin{longtable}[]{@{}
  >{\raggedright\arraybackslash}p{(\linewidth - 2\tabcolsep) * \real{0.2143}}
  >{\raggedright\arraybackslash}p{(\linewidth - 2\tabcolsep) * \real{0.7857}}@{}}
\toprule\noalign{}
\begin{minipage}[b]{\linewidth}\raggedright
Goal
\end{minipage} & \begin{minipage}[b]{\linewidth}\raggedright
Normative consequence
\end{minipage} \\
\midrule\noalign{}
\endhead
\bottomrule\noalign{}
\endlastfoot
Syntactic familiarity & The surface grammar follows
JavaScript/TypeScript conventions unless this document states
otherwise. \\
Paradigm flexibility & Procedural, object-oriented, and functional
constructs are first-class and freely composable. \\
Optional typing & Dynamic typing is the default; strict typing is opt-in
per module, function, or declaration. \\
Triple execution & The interpreter, bytecode VM, and AOT compiler MUST
preserve the same observable semantics. \\
Native concurrency & Threads, mutexes, and asynchronous functions are
language primitives, not library add-ons. \\
Self-hosting & The standard library MUST be written in pure Dryad and
MUST use only the specified intrinsics. \\
Zero surprises & Every construct in this document has exactly one
defined behaviour; undefined behaviour is not permitted. \\
\end{longtable}

\subsection{Reading guide}\label{reading-guide}

A reader implementing an interpreter SHOULD read Clauses 1--15 in order.
A reader implementing the bytecode VM SHOULD additionally read Clause
19. A reader implementing the AOT compiler SHOULD read Clauses 19 and
20. A reader implementing the standard library SHOULD read Clauses 8,
14, 17, and 18.

\section{1 Scope}\label{scope}

\textbf{1.1} This document specifies the syntax, static semantics, and
runtime semantics of the Dryad programming language, together with the
normative requirements for its execution engines, native runtime
interface, standard library, bytecode format, and ahead-of-time
compilation pipeline.

\textbf{1.2} This document specifies the observable behaviour of a
conforming implementation. Observable behaviour comprises:

\begin{enumerate}
\def\labelenumi{\arabic{enumi}.}
\tightlist
\item
  the result of evaluating every well-formed program;
\item
  the sequence of side effects produced during evaluation;
\item
  the errors and diagnostics reported for ill-formed programs;
\item
  the values returned by standard-library functions; and
\item
  the concurrency guarantees defined in Clause 15.
\end{enumerate}

\textbf{1.3} This document does not specify:

\begin{enumerate}
\def\labelenumi{\arabic{enumi}.}
\tightlist
\item
  the hardware, operating system, or host environment on which an
  implementation runs, except where a boundary is explicitly defined;
\item
  the internal data structures, algorithms, or source-code organisation
  of an implementation;
\item
  user interfaces, editors, debuggers, package registries, or build
  tools, except where their behaviour is observable through a normative
  interface;
\item
  the contents of third-party libraries; or
\item
  performance, except where a normative requirement is explicitly
  stated.
\end{enumerate}

\textbf{1.4} A program is \emph{well-formed} if it satisfies every early
error rule in this document. A well-formed program has exactly one
specified behaviour. This document does not define undefined behaviour;
where a program violates a precondition, the implementation MUST report
a runtime error as specified in Clause 14.

\textbf{1.5} The unit of conformance is a \emph{conforming
implementation}, as defined in Clause 2. Conformance is claimed per
conformance class, not per feature.

\textbf{1.6} This document is organised as follows:

\begin{longtable}[]{@{}
  >{\raggedright\arraybackslash}p{(\linewidth - 2\tabcolsep) * \real{0.5000}}
  >{\raggedright\arraybackslash}p{(\linewidth - 2\tabcolsep) * \real{0.5000}}@{}}
\toprule\noalign{}
\begin{minipage}[b]{\linewidth}\raggedright
Clauses
\end{minipage} & \begin{minipage}[b]{\linewidth}\raggedright
Subject
\end{minipage} \\
\midrule\noalign{}
\endhead
\bottomrule\noalign{}
\endlastfoot
1--6 & Scope, conformance, references, terminology, notation,
overview \\
7 & Lexical grammar \\
8 & Types and values \\
9 & Expressions \\
10 & Statements \\
11 & Functions and classes \\
12 & Modules \\
13 & Execution model \\
14 & Error system \\
15 & Concurrency model \\
16 & Foreign function interface \\
17 & Runtime architecture and intrinsics \\
18 & Standard library \\
19 & Bytecode and intermediate representation \\
20 & Ahead-of-time compilation \\
21--24 & Paradigms, formal grammar, conformance testing, limitations \\
Annexes A--G & Registries, editorial decisions, proposals,
bibliography \\
\end{longtable}

\section{2 Conformance}\label{conformance}

\textbf{2.1} A \emph{conforming implementation} is a program that
implements every normative requirement of this document for at least one
conformance class, and that does not contradict any other normative
requirement.

\textbf{2.2} This document defines four conformance classes.

\begin{longtable}[]{@{}
  >{\raggedright\arraybackslash}p{(\linewidth - 6\tabcolsep) * \real{0.1429}}
  >{\raggedright\arraybackslash}p{(\linewidth - 6\tabcolsep) * \real{0.2245}}
  >{\raggedright\arraybackslash}p{(\linewidth - 6\tabcolsep) * \real{0.3673}}
  >{\raggedright\arraybackslash}p{(\linewidth - 6\tabcolsep) * \real{0.2653}}@{}}
\toprule\noalign{}
\begin{minipage}[b]{\linewidth}\raggedright
Class
\end{minipage} & \begin{minipage}[b]{\linewidth}\raggedright
Identifier
\end{minipage} & \begin{minipage}[b]{\linewidth}\raggedright
Required clauses
\end{minipage} & \begin{minipage}[b]{\linewidth}\raggedright
Description
\end{minipage} \\
\midrule\noalign{}
\endhead
\bottomrule\noalign{}
\endlastfoot
Core & \texttt{core-1.0} & 7--18 & Lexical analysis, parsing, static
semantics, interpretation, modules, errors, concurrency, FFI, runtime
interface, standard library. \\
Virtual Machine & \texttt{vm-1.0} & 7--19 & Core, plus the bytecode
compiler, bytecode verifier, and bytecode VM. \\
Ahead-of-Time & \texttt{aot-1.0} & 7--20 & Core, plus the bytecode
pipeline and the AOT compiler. \\
Full & \texttt{full-1.0} & 7--20 & Core, VM, and AOT in a single
implementation. \\
\end{longtable}

\textbf{2.3} An implementation that claims a conformance class MUST:

\begin{enumerate}
\def\labelenumi{\arabic{enumi}.}
\tightlist
\item
  accept every program that this document classifies as well-formed;
\item
  reject every program that this document classifies as ill-formed, with
  a diagnostic as specified in Clause 14;
\item
  produce the specified result for every well-formed program;
\item
  preserve the concurrency guarantees of Clause 15; and
\item
  satisfy every MUST or MUST NOT requirement in the clauses required by
  its class.
\end{enumerate}

\textbf{2.4} An implementation MAY provide extensions provided that all
of the following hold:

\begin{enumerate}
\def\labelenumi{\arabic{enumi}.}
\tightlist
\item
  the extension is not observable in a program that uses only the
  constructs specified in this document;
\item
  the extension does not alter the meaning of any specified construct;
\item
  the implementation documents the extension and the conditions under
  which it is active; and
\item
  disabling the extension restores a conforming implementation.
\end{enumerate}

An implementation MUST NOT use the availability of an extension to
justify a deviation from a normative requirement.

\textbf{2.5} Certain behaviours are designated
\emph{implementation-defined} or \emph{implementation-dependent}. For
implementation-defined behaviour, a conforming implementation MUST
document the chosen behaviour. For implementation-dependent behaviour,
the implementation may choose freely and is not required to document the
choice. This document identifies such behaviours explicitly; no other
behaviour is implementation-defined or implementation-dependent.

\textbf{2.6} A conforming implementation SHOULD pass the conformance
test suite maintained by the Dryad Language Project for its claimed
class. Passing the test suite is evidence of conformance but is not a
substitute for the normative requirements of this document.

\textbf{2.7} An implementation MUST report its edition and conformance
class when invoked with \texttt{-\/-version} or an equivalent mechanism.
The report MUST identify this document as \texttt{DRYAD-SPEC-1.0-DRAFT}.

\textbf{2.8} Conformance is not affected by:

\begin{enumerate}
\def\labelenumi{\arabic{enumi}.}
\tightlist
\item
  the execution tier used to evaluate a program, provided the observable
  behaviour is preserved;
\item
  the host operating system, provided the specified platform
  abstractions are implemented;
\item
  the presence or absence of optional optimisations; or
\item
  the textual form of a diagnostic, provided it contains the information
  required by Clause 14.
\end{enumerate}

\textbf{2.9} An implementation that cannot support a required facility
on a given host MUST reject programs that require that facility at the
earliest specified point, with an error of category \emph{unsupported}
(Clause 14). Silent degradation is not conforming.

\section{3 Normative References}\label{normative-references}

\textbf{3.1} The following documents are indispensable for the
application of this specification. For dated references, only the cited
edition applies. For undated references, the latest edition applies.

\begin{longtable}[]{@{}
  >{\raggedright\arraybackslash}p{(\linewidth - 4\tabcolsep) * \real{0.2750}}
  >{\raggedright\arraybackslash}p{(\linewidth - 4\tabcolsep) * \real{0.1750}}
  >{\raggedright\arraybackslash}p{(\linewidth - 4\tabcolsep) * \real{0.5500}}@{}}
\toprule\noalign{}
\begin{minipage}[b]{\linewidth}\raggedright
Reference
\end{minipage} & \begin{minipage}[b]{\linewidth}\raggedright
Title
\end{minipage} & \begin{minipage}[b]{\linewidth}\raggedright
Use in this document
\end{minipage} \\
\midrule\noalign{}
\endhead
\bottomrule\noalign{}
\endlastfoot
RFC 2119 & Key words for use in RFCs to Indicate Requirement Levels &
Normative keywords (Clause 5.1) \\
RFC 8174 & Ambiguity of Uppercase vs Lowercase in RFC 2119 Key Words &
Normative keyword interpretation \\
RFC 3629 & UTF-8, a transformation format of ISO 10646 & Source-text
encoding (Clause 7.1) \\
ISO/IEC 10646 & Information technology --- Universal Coded Character Set
(UCS) & Character repertoire \\
The Unicode Standard & Version 15.1 or later & Character properties,
identifier classes \\
IEEE 754-2019 & IEEE Standard for Floating-Point Arithmetic & Number
semantics (Clause 8.2.3) \\
ISO/IEC 9899:2018 & Programming languages --- C & Foreign function
interface ABI (Clause 16) \\
System V ABI & System V Application Binary Interface, x86-64 and ARM64
editions & Calling conventions (Clauses 16, 20) \\
ELF gABI & System V gABI, Object File Format & AOT object output (Clause
20) \\
PE/COFF & Microsoft Portable Executable and Common Object File Format &
AOT object output (Clause 20) \\
ECMA-262 & ECMAScript Language Specification & Informative reference for
algorithm style and grammar notation \\
\end{longtable}

\textbf{3.2} The following documents are informative references. They
describe the design context of this specification but do not create
conformance requirements.

\begin{longtable}[]{@{}ll@{}}
\toprule\noalign{}
Reference & Title \\
\midrule\noalign{}
\endhead
\bottomrule\noalign{}
\endlastfoot
DRYAD-TF-2.0 & Dryad Theoretical Foundation, version 2.0 \\
DRYAD-CONV-1.0 & Dryad Conventions, version 1.0 (superseded) \\
LLVM Language Reference Manual & LLVM IR used by the AOT pipeline \\
OpenSSL documentation & TLS backend named by the standard library \\
POSIX.1-2017 & Portable Operating System Interface \\
\end{longtable}

\textbf{3.3} Where a normative reference is updated, the Dryad Language
Project MAY publish an erratum adopting the update. Until such an
erratum is published, the dated edition cited in Clause 3.1 applies.

\section{4 Terms and Definitions}\label{terms-and-definitions}

For the purposes of this document, the following terms and definitions
apply. Terms are listed alphabetically. Cross-references use clause
numbers.

\textbf{Abstract operation.} A named algorithm defined by this document.
Abstract operations are not callable from Dryad source text; they define
semantics.

\textbf{Agent.} An independent unit of execution with its own execution
context and job queue. A Dryad program executes in at least one agent.
Threads created by \texttt{thread\ function} execute in separate agents
(Clause 15).

\textbf{AOT compilation.} Ahead-of-time compilation: translation of a
Dryad program to native code before execution, as specified in Clause
20.

\textbf{Binding.} The association between an identifier and a value
within an environment record.

\textbf{Buffer.} A contiguous, bounds-checked byte region with an
explicit length and capacity (Clause 17.5).

\textbf{Bytecode.} The stack-based instruction format defined in Clause
19.

\textbf{Call frame.} The runtime record that stores the state of an
active function invocation.

\textbf{Class.} A value that defines a nominal object type, its
constructor, its instance members, and its static members (Clause 11).

\textbf{Completion record.} A specification value that records how
evaluation of a statement or expression terminated. A completion record
is one of \emph{normal}, \emph{break}, \emph{continue}, \emph{return},
or \emph{throw} (Clause 10.2).

\textbf{Conforming implementation.} An implementation that satisfies
Clause 2.

\textbf{Directive.} A syntactic construct beginning with \texttt{\#}
that instructs the implementation, such as a native-module directive
\texttt{\#\textless{}io\textgreater{}} (Clause 7.10), or a
string-literal directive such as \texttt{"use\ strict\ types"} (Clause
8.6).

\textbf{Dynamic type.} The type of a value at runtime. Every value has
exactly one dynamic type (Clause 8).

\textbf{Early error.} A condition that makes a program ill-formed and
that MUST be detected before execution, during parsing or static
analysis (Clause 14.8).

\textbf{Environment record.} A specification data structure that maps
identifiers to bindings and that may have an outer environment reference
(Clause 13.2).

\textbf{Execution context.} The specification state associated with the
evaluation of a program, module, function, or block. An execution
context includes a lexical environment, a variable environment, and a
\texttt{this} binding (Clause 13.2).

\textbf{Export.} A binding made available by a module to importing
modules (Clause 12.3).

\textbf{Expression.} A syntactic construct that evaluates to a value
(Clause 9).

\textbf{FFI.} Foreign function interface: the boundary between Dryad and
code implemented in another language (Clause 16).

\textbf{Function.} A callable value with a parameter list, a body, and a
captured lexical environment (Clause 11).

\textbf{HAL.} Hardware abstraction layer: the small set of
platform-specific operations beneath the intrinsic layer (Clause 17.4).

\textbf{Host.} The environment in which a conforming implementation
executes, including its operating system, filesystem, network stack, and
clock.

\textbf{Identifier.} A lexical token that names a binding, property,
label, or module (Clause 7.5).

\textbf{Implementation-defined.} A behaviour for which a conforming
implementation MUST document its choice (Clause 2.5).

\textbf{Implementation-dependent.} A behaviour for which a conforming
implementation may choose freely without documentation (Clause 2.5).

\textbf{Intrinsic.} A primitive operation exposed by the runtime and
declared in Dryad with \texttt{@intrinsic} (Clause 17.2).

\textbf{Job.} A unit of work scheduled on an agent's job queue, used to
implement promises and asynchronous functions (Clause 15.4).

\textbf{Lexical grammar.} The grammar that defines how source text is
divided into tokens (Clause 7).

\textbf{Lexical token.} An indivisible lexical unit produced by the
lexer (Clause 7.4).

\textbf{Module.} A source-text unit with its own scope, exports, and
module record (Clause 12).

\textbf{Module record.} The specification data structure that records a
module's exports, imports, and evaluation state (Clause 12.1).

\textbf{Module specifier.} The string or identifier that names a module
in an import declaration (Clause 12.4).

\textbf{Native directive.} The token
\texttt{\#\textless{}identifier\textgreater{}} that loads a native
module (Clause 7.10).

\textbf{Native module.} A host-implemented module exposed through the
FFI or the intrinsic layer (Clauses 16, 17).

\textbf{Native type.} A machine-level type used by the runtime, AOT
compiler, and FFI, such as \texttt{f64}, \texttt{i32}, or \texttt{usize}
(Clause 8.3).

\textbf{Normal completion.} A completion record indicating that
evaluation continued to the next construct without an abrupt transfer of
control (Clause 10.2).

\textbf{Object.} A value that groups properties and has an identity
independent of its contents (Clause 8.3).

\textbf{Observable behaviour.} The behaviour listed in Clause 1.2.

\textbf{Program.} A collection of one or more modules, one of which is
the entry module. Execution begins at the entry module (Clause 13.7).

\textbf{Promise.} A value representing the eventual result of an
asynchronous operation (Clause 15.3).

\textbf{Property.} A named or indexed member of an object.

\textbf{Realm.} The set of intrinsic objects and global bindings
associated with an agent. A realm is created when an agent is created
(Clause 13.1).

\textbf{Scope.} The region of source text in which a binding is visible,
together with the rules that determine visibility (Clause 13.3).

\textbf{Shape.} A compiler-internal description of an object's property
layout, used by the AOT compiler (Clause 20.4).

\textbf{Source text.} A sequence of Unicode code points encoded in UTF-8
(Clause 7.1).

\textbf{Standard library (STDLIB).} The collection of modules specified
in Clause 18 and implemented in pure Dryad.

\textbf{Statement.} A syntactic construct that performs an action and
produces a completion record rather than a value (Clause 10).

\textbf{Static semantics.} Rules that constrain well-formed programs and
that are enforced before execution, including early errors and type
checking (Clauses 8, 22).

\textbf{Strict mode.} The optional static typing mode enabled by
\texttt{"use\ strict\ types"} (Clause 8.6).

\textbf{Syscall.} A primitive operation in the intrinsic set that
corresponds to an operating-system service (Clause 17.3).

\textbf{Thread.} An independent agent created by a
\texttt{thread\ function} declaration or a thread-creation operation
(Clause 15.1).

\textbf{Type.} A set of values. The dynamic type of a value determines
the operations that are defined on it (Clause 8).

\textbf{Type annotation.} A syntactic type expression attached to a
declaration, parameter, return position, or expression (Clause 8.4).

\textbf{Value.} A member of a Dryad type. Values are the results of
expressions and the contents of bindings (Clause 8).

\textbf{Variable.} A binding whose value may be changed after
declaration.

\textbf{Well-formed program.} A program that satisfies every early error
rule of this document (Clause 1.4).

\section{5 Notational Conventions}\label{notational-conventions}

\subsection{5.1 Normative keywords}\label{normative-keywords}

\textbf{5.1.1} The key words MUST, MUST NOT, REQUIRED, SHALL, SHALL NOT,
SHOULD, SHOULD NOT, RECOMMENDED, NOT RECOMMENDED, MAY, and OPTIONAL in
this document are to be interpreted as described in RFC 2119 and RFC
8174.

\textbf{5.1.2} These keywords are normative only when they appear in
uppercase. Lowercase uses of the same words are descriptive.

\textbf{5.1.3} A requirement stated without a keyword is normative if it
appears in a normative clause and describes observable behaviour, a
well-formedness condition, or a conformance requirement.

\textbf{5.1.4} The word \emph{can} describes possibility; the word
\emph{cannot} describes impossibility under this specification. Neither
word creates an optional implementation choice.

\subsection{5.2 Grammar notation}\label{grammar-notation}

\textbf{5.2.1} Grammars in this document use an extended Backus--Naur
form (EBNF) with the following metacharacters.

\begin{longtable}[]{@{}
  >{\raggedright\arraybackslash}p{(\linewidth - 2\tabcolsep) * \real{0.5263}}
  >{\raggedright\arraybackslash}p{(\linewidth - 2\tabcolsep) * \real{0.4737}}@{}}
\toprule\noalign{}
\begin{minipage}[b]{\linewidth}\raggedright
Notation
\end{minipage} & \begin{minipage}[b]{\linewidth}\raggedright
Meaning
\end{minipage} \\
\midrule\noalign{}
\endhead
\bottomrule\noalign{}
\endlastfoot
\texttt{::=} & Production definition. \\
\texttt{\textless{}name\textgreater{}} & Nonterminal symbol
\texttt{name}. \\
\texttt{"text"} & Terminal symbol consisting of the exact characters
\texttt{text}. \\
\texttt{A\ \textbackslash{}\textbar{}\ B} & Alternative: either
\texttt{A} or \texttt{B}. \\
\texttt{{[}\ A\ {]}} & Optional occurrence of \texttt{A} (zero or
one). \\
\texttt{\{\ A\ \}} & Repetition of \texttt{A} (zero or more). \\
\texttt{(\ A\ )} & Grouping. \\
\texttt{ε} & The empty production. \\
\texttt{(*\ text\ *)} & Informative comment inside a grammar. \\
\end{longtable}

\textbf{5.2.2} A production of the form
\texttt{\textless{}A\textgreater{}\ ::=\ \textless{}B\textgreater{}\ \{\ ","\ \textless{}B\textgreater{}\ \}}
denotes a possibly empty comma-separated sequence of
\texttt{\textless{}B\textgreater{}}. A trailing comma is permitted only
where a production says so explicitly.

\textbf{5.2.3} Lexical grammar productions are identified by the prefix
\texttt{Lexical} and are defined in Clause 7. Syntactic grammar
productions are defined in Clause 22 and referenced throughout this
document. When a syntactic production references a token, the token name
is written in angle brackets, for example
\texttt{\textless{}Identifier\textgreater{}}.

\textbf{5.2.4} The lexical grammar is applied before the syntactic
grammar. The lexer consumes the longest sequence of input characters
that forms a valid token, subject to the rules of Clause 7.11.
Whitespace and comments separate tokens but are not tokens themselves.

\textbf{5.2.5} A syntactic grammar is unambiguous if every well-formed
token sequence has exactly one parse tree. This document defines an
unambiguous grammar. An implementation MUST reject a token sequence for
which no parse tree exists and MUST report an early error if more than
one parse tree exists.

\textbf{5.2.6} Grammar productions do not by themselves impose
evaluation order. Evaluation order is defined by the algorithms in
Clauses 9 through 13.

\subsection{5.3 Algorithm conventions}\label{algorithm-conventions}

\textbf{5.3.1} Semantic algorithms are written as numbered steps. Steps
are performed in order unless a step transfers control.

\textbf{5.3.2} The phrase \emph{Let x be \ldots{}} introduces a local
specification variable \texttt{x}. Specification variables are distinct
from Dryad variables and are not observable by a program.

\textbf{5.3.3} The phrase \emph{If condition, then \ldots{}} executes
the indented steps when \texttt{condition} is true. \emph{Otherwise}
executes them when \texttt{condition} is false. \emph{If condition, then
\ldots; otherwise \ldots{}} is equivalent.

\textbf{5.3.4} The phrase \emph{Return value} terminates the abstract
operation and yields \texttt{value}. The phrase \emph{Return} without a
value yields the normal completion \texttt{Normal(empty)}.

\textbf{5.3.5} The phrase \emph{Assert: condition} states an invariant
that MUST hold whenever the algorithm reaches that step. A violation
indicates an error in this specification or in a conforming
implementation, not in the source program.

\textbf{5.3.6} An abstract operation whose name ends in \texttt{?} MAY
return an abrupt completion. In an algorithm, the notation
\texttt{result\ =\ Operation?} means: perform \texttt{Operation}; if it
returns an abrupt completion, return that completion immediately;
otherwise, assign its value to \texttt{result}.

\textbf{5.3.7} Abstract operations named in CamelCase, such as
\texttt{GetValue} or \texttt{ResolveBinding}, are specification
operations. They are not callable from Dryad source text.

\textbf{5.3.8} An algorithm MAY state that behaviour is
\emph{implementation-defined} or \emph{implementation-dependent}. Those
terms have the meanings given in Clause 2.5.

\subsection{5.4 Completion records}\label{completion-records}

\textbf{5.4.1} A completion record is a specification value of one of
the following forms.

\begin{longtable}[]{@{}
  >{\raggedright\arraybackslash}p{(\linewidth - 4\tabcolsep) * \real{0.2500}}
  >{\raggedright\arraybackslash}p{(\linewidth - 4\tabcolsep) * \real{0.3750}}
  >{\raggedright\arraybackslash}p{(\linewidth - 4\tabcolsep) * \real{0.3750}}@{}}
\toprule\noalign{}
\begin{minipage}[b]{\linewidth}\raggedright
Form
\end{minipage} & \begin{minipage}[b]{\linewidth}\raggedright
Payload
\end{minipage} & \begin{minipage}[b]{\linewidth}\raggedright
Meaning
\end{minipage} \\
\midrule\noalign{}
\endhead
\bottomrule\noalign{}
\endlastfoot
\texttt{Normal(value)} & A value or \texttt{empty} & Evaluation
continued normally. \\
\texttt{Break(label)} & A label or \texttt{empty} & Control transfers to
the matching labelled statement. \\
\texttt{Continue(label)} & A label or \texttt{empty} & Control transfers
to the next iteration of the matching loop. \\
\texttt{Return(value)} & A value & Control returns from the current
function. \\
\texttt{Throw(value)} & A value & An exception is propagating. \\
\end{longtable}

\textbf{5.4.2} A completion record is \emph{abrupt} if it is not
\texttt{Normal}. A completion record is \emph{empty} if it carries no
value.

\textbf{5.4.3} Statement evaluation in Clause 10 yields completion
records. Expression evaluation in Clause 9 yields values or throws; an
expression never yields \texttt{Break}, \texttt{Continue}, or
\texttt{Return}.

\subsection{5.5 Type notation}\label{type-notation}

\textbf{5.5.1} This document uses the following notation for value
types.

\begin{longtable}[]{@{}
  >{\raggedright\arraybackslash}p{(\linewidth - 2\tabcolsep) * \real{0.5263}}
  >{\raggedright\arraybackslash}p{(\linewidth - 2\tabcolsep) * \real{0.4737}}@{}}
\toprule\noalign{}
\begin{minipage}[b]{\linewidth}\raggedright
Notation
\end{minipage} & \begin{minipage}[b]{\linewidth}\raggedright
Meaning
\end{minipage} \\
\midrule\noalign{}
\endhead
\bottomrule\noalign{}
\endlastfoot
\texttt{Null} & The type whose only value is \texttt{null}. \\
\texttt{Boolean} & \texttt{true} and \texttt{false}. \\
\texttt{Number} & IEEE 754 binary64 values, including \texttt{NaN} and
infinities. \\
\texttt{String} & Finite sequences of Unicode scalar values. \\
\texttt{Object} & Objects with identity, including arrays, tuples,
functions, and class instances. \\
\texttt{Array\textless{}T\textgreater{}} & An ordered, dynamically sized
sequence of \texttt{T}. \\
\texttt{Tuple\textless{}T1,\ ...,\ Tn\textgreater{}} & A fixed-length
sequence with element types \texttt{T1} through \texttt{Tn}. \\
\texttt{Function\textless{}P...,\ R\textgreater{}} & A callable value
with parameters \texttt{P...} and result \texttt{R}. \\
\texttt{Class\textless{}C\textgreater{}} & The type of instances of
class \texttt{C}. \\
\texttt{Buffer} & The byte-buffer type of Clause 17.5. \\
\texttt{Promise\textless{}T\textgreater{}} & The asynchronous-result
type of Clause 15.3. \\
\texttt{Result\textless{}T,\ E\textgreater{}} & A value that is either
\texttt{Ok(T)} or \texttt{Err(E)} (Clause 8.3.7). \\
\texttt{Thread\textless{}T\textgreater{}} & A handle to a thread
producing \texttt{T} (Clause 15.1). \\
\texttt{Mutex} & A mutual-exclusion primitive (Clause 15.2). \\
\texttt{any} & The union of all value types. \\
\texttt{never} & The empty type; an expression of type \texttt{never}
does not return normally. \\
\texttt{T?} & Shorthand for \texttt{T\ \textbar{}\ Null}. \\
\texttt{T\ \textbar{}\ U} & The union of \texttt{T} and \texttt{U}. \\
\end{longtable}

\textbf{5.5.2} This document uses the following notation for native
types, which are visible to the AOT compiler, the FFI, and intrinsic
declarations.

\begin{longtable}[]{@{}ll@{}}
\toprule\noalign{}
Notation & Meaning \\
\midrule\noalign{}
\endhead
\bottomrule\noalign{}
\endlastfoot
\texttt{f64} & IEEE 754 binary64. \\
\texttt{i8}, \texttt{i16}, \texttt{i32}, \texttt{i64} & Signed
two's-complement integers. \\
\texttt{u8}, \texttt{u16}, \texttt{u32}, \texttt{u64} & Unsigned
integers. \\
\texttt{usize}, \texttt{isize} & Pointer-sized unsigned and signed
integers. \\
\texttt{bool} & A native boolean. \\
\texttt{ptr\textless{}T\textgreater{}} & A raw pointer to \texttt{T}. \\
\texttt{void} & No value. \\
\end{longtable}

\textbf{5.5.3} \texttt{T\ \textless{}:\ U} means that \texttt{T} is a
subtype of \texttt{U}; every value of type \texttt{T} is also a value of
type \texttt{U}. \texttt{T\ \textasciitilde{}\ U} means that \texttt{T}
and \texttt{U} are consistent: their intersection is non-empty and
values may flow between them under the rules of Clause 8.5.

\textbf{5.5.4} \texttt{typeof(v)} denotes the dynamic type of value
\texttt{v}. \texttt{IsCallable(v)} denotes whether \texttt{v} is a
function or class constructor. \texttt{IsObject(v)} denotes whether
\texttt{v} is an object.

\subsection{5.6 Mathematical notation}\label{mathematical-notation}

\textbf{5.6.1} \texttt{R} denotes the real numbers, \texttt{Z} the
integers, and \texttt{B} the booleans.

\textbf{5.6.2} A sequence is written
\texttt{\textless{}a1,\ ...,\ an\textgreater{}}; an empty sequence is
\texttt{\textless{}\textgreater{}}. The length of sequence \texttt{s} is
\texttt{len(s)}. Indexing is zero-based and written \texttt{s{[}i{]}}.

\textbf{5.6.3} A map is written
\texttt{\{k1\ -\textgreater{}\ v1,\ ...,\ kn\ -\textgreater{}\ vn\}};
\texttt{m{[}k{]}} denotes lookup and
\texttt{m{[}k\ -\textgreater{}\ v{]}} denotes extension or update.

\textbf{5.6.4} The interval \texttt{{[}a,\ b)} denotes the set of
integers \texttt{x} such that \texttt{a\ ≤\ x\ \textless{}\ b}.
\texttt{{[}a,\ b{]}} includes \texttt{b}.

\textbf{5.6.5} \texttt{floor(x)} denotes the largest integer less than
or equal to \texttt{x}; \texttt{trunc(x)} denotes the integer part of
\texttt{x} toward zero.

\textbf{5.6.6} \texttt{SameValue(x,\ y)} denotes the specification
equality relation of Clause 8.9. \texttt{SameValueZero(x,\ y)} is
\texttt{SameValue} except that it treats \texttt{+0} and \texttt{-0} as
equal.

\subsection{5.7 Normative and informative
marking}\label{normative-and-informative-marking}

\textbf{5.7.1} A clause is normative unless its heading or introductory
paragraph marks it informative. Annexes A through D are normative;
Annexes E through G are informative.

\textbf{5.7.2} Within a normative clause, a paragraph beginning with
\textbf{NOTE} is informative and does not create a requirement.

\textbf{5.7.3} Examples are informative unless the example is introduced
as a normative test vector.

\subsection{5.8 Document versioning}\label{document-versioning}

\textbf{5.8.1} This document is Edition 1.0, Draft. Clause numbers are
stable within the edition. An erratum MUST NOT renumber a clause; it MAY
add subclauses with a new final component.

\textbf{5.8.2} A future edition MAY renumber clauses only if it
publishes a migration table mapping old clause numbers to new ones.

\section{6 Overview}\label{overview}

\subsection{6.1 Language identity}\label{language-identity}

\textbf{6.1.1} Dryad is a statically scoped, dynamically typed,
multi-paradigm programming language. It supports procedural,
object-oriented, and functional programming in a single surface syntax.

\textbf{6.1.2} Dryad source text is compiled to an internal program
representation and may be executed by three interchangeable tiers: a
tree-walking interpreter, a bytecode virtual machine, or an
ahead-of-time compiled native binary.

\textbf{6.1.3} The language provides first-class functions, closures,
classes with single inheritance, interfaces, modules, exceptions,
asynchronous functions, operating-system threads, mutexes, a foreign
function interface, and a standard library written entirely in Dryad.

\subsection{6.2 Design axioms}\label{design-axioms}

\emph{This subclause is informative. The axioms guide interpretation
where this document leaves room for judgement; they do not override an
explicit requirement.}

\textbf{Axiom 1 --- Clarity over performance.} Readability and
predictable behaviour take precedence over premature optimisation. An
optimisation is conforming only if it preserves observable behaviour.

\textbf{Axiom 2 --- Syntactic consistency.} Syntactic constructs follow
uniform, predictable patterns. Similar operations use similar notation.

\textbf{Axiom 3 --- Informative errors.} Every error reports a category,
a numeric code, a source location, and, where the cause is determinable,
a suggested correction (Clause 14).

\textbf{Axiom 4 --- Composability.} Language primitives combine without
special cases. A value of a given type may be used wherever that type is
accepted.

\textbf{Axiom 5 --- Zero surprises.} Every construct has exactly one
specified behaviour. The implementation MUST NOT depend on unspecified
state, ambient configuration, or evaluation order that this document
leaves undefined.

\subsection{6.3 Surface language}\label{surface-language}

\textbf{6.3.1} A minimal Dryad program is a sequence of statements. The
following program is well-formed and prints \texttt{42}:

\begin{Shaded}
\begin{Highlighting}[]
\NormalTok{function answer() \{}
\NormalTok{    return 42;}
\NormalTok{\}}

\NormalTok{print(answer());}
\end{Highlighting}
\end{Shaded}

\textbf{6.3.2} The principal syntactic categories are:

\begin{longtable}[]{@{}
  >{\raggedright\arraybackslash}p{(\linewidth - 4\tabcolsep) * \real{0.3571}}
  >{\raggedright\arraybackslash}p{(\linewidth - 4\tabcolsep) * \real{0.3571}}
  >{\raggedright\arraybackslash}p{(\linewidth - 4\tabcolsep) * \real{0.2857}}@{}}
\toprule\noalign{}
\begin{minipage}[b]{\linewidth}\raggedright
Category
\end{minipage} & \begin{minipage}[b]{\linewidth}\raggedright
Examples
\end{minipage} & \begin{minipage}[b]{\linewidth}\raggedright
Clause
\end{minipage} \\
\midrule\noalign{}
\endhead
\bottomrule\noalign{}
\endlastfoot
Declarations & \texttt{let}, \texttt{const}, \texttt{function},
\texttt{class}, \texttt{import}, \texttt{export} & 10--12 \\
Statements & \texttt{if}, \texttt{while}, \texttt{for}, \texttt{try},
\texttt{return}, \texttt{throw} & 10 \\
Expressions & operators, calls, member access, \texttt{match}, arrow
functions & 9 \\
Types & primitives, arrays, tuples, functions, classes, native types &
8 \\
Modules & named, namespace, side-effect, and identifier imports & 12 \\
Directives & \texttt{\#\textless{}io\textgreater{}},
\texttt{"use\ strict\ types"}, \texttt{@compile(...)} & 7, 8, 13 \\
\end{longtable}

\subsection{6.4 Execution architecture}\label{execution-architecture}

\textbf{6.4.1} A conforming implementation MUST provide at least one
execution tier and SHOULD provide all three. The tiers are:

\begin{longtable}[]{@{}
  >{\raggedright\arraybackslash}p{(\linewidth - 4\tabcolsep) * \real{0.2308}}
  >{\raggedright\arraybackslash}p{(\linewidth - 4\tabcolsep) * \real{0.2692}}
  >{\raggedright\arraybackslash}p{(\linewidth - 4\tabcolsep) * \real{0.5000}}@{}}
\toprule\noalign{}
\begin{minipage}[b]{\linewidth}\raggedright
Tier
\end{minipage} & \begin{minipage}[b]{\linewidth}\raggedright
Input
\end{minipage} & \begin{minipage}[b]{\linewidth}\raggedright
Requirement
\end{minipage} \\
\midrule\noalign{}
\endhead
\bottomrule\noalign{}
\endlastfoot
Interpreter & Abstract syntax tree & Evaluate the tree directly. \\
Bytecode VM & Bytecode module (Clause 19) & Execute verified
bytecode. \\
AOT & Native object code (Clause 20) & Execute ahead-of-time compiled
code. \\
\end{longtable}

\textbf{6.4.2} All tiers MUST be observationally equivalent for every
well-formed program. Tier selection is not observable except through the
timing and the diagnostic facilities explicitly defined in Clause 20.

\textbf{6.4.3} The interpreter MUST collect the runtime type
observations required by the metadata header of Clause 19.8 when the
implementation also provides a bytecode or AOT tier. An implementation
that provides only the interpreter is not required to emit metadata.

\textbf{6.4.4} The \texttt{@compile} decorator (Clause 13.9) selects the
execution tier for a function or block. In the absence of a decorator,
the tier is selected by the implementation or by the implementation's
command-line interface.

\subsection{6.5 Native boundary and
self-hosting}\label{native-boundary-and-self-hosting}

\textbf{6.5.1} A Dryad implementation consists of a native runtime and a
standard library. The native runtime MUST expose only the intrinsics
specified in Clause 17.3 and Annex C. The standard library MUST be
implemented in pure Dryad and MUST obtain all host services through
those intrinsics.

\textbf{6.5.2} An implementation MAY provide additional host facilities
through the foreign function interface of Clause 16. Such facilities are
not part of the standard library and do not extend the intrinsic set.

\textbf{6.5.3} The native runtime MUST NOT contain language-specific
standard library logic beyond the intrinsics. In particular, collection
types, formatting, HTTP, JSON, cryptography, and path handling are
implemented in Dryad.

\subsection{6.6 Document organisation}\label{document-organisation}

\textbf{6.6.1} The normative clauses are organised by language layer.
Clauses 7--12 define the source language: lexical structure, values and
types, expressions, statements, declarations, and modules. Clause 13
defines how those constructs execute. Clause 14 defines errors. Clause
15 defines concurrency. Clauses 16 and 17 define the native boundary.
Clause 18 defines the library. Clauses 19 and 20 define the compiled
representations.

\textbf{6.6.2} Clause 22 contains the complete grammar. A production in
Clause 22 is normative. Syntax shown in other clauses is explanatory and
MUST agree with Clause 22; if it does not, Clause 22 prevails and the
discrepancy is an erratum.

\textbf{6.6.3} Annex A is the error-code registry. Annex B is the opcode
registry. Annex C is the intrinsic registry. Annex D is the
standard-library index. Annex E records editorial decisions against the
conventions draft. Annex F is the informative proposals annex. Annex G
is the bibliography.

\end{document}
